\documentclass[12pt]{article}
\textwidth 7.0in
\oddsidemargin -0.4in
\evensidemargin -0.4in
\textheight 9.0in 
\pagestyle{empty}
\usepackage{amsmath, amssymb, amscd, amsfonts}
\usepackage{amsthm}
\usepackage{latexsym, epsfig, color, verbatim, xcolor, color}
\usepackage{url}


\newcommand{\Sym}{\operatorname{Sym}}
\newcommand{\GL}{\operatorname{GL}}
\newcommand{\SL}{\operatorname{SL}}
\newcommand{\leg}[2]{\genfrac{(}{)}{}{}{#1}{#2}}
\newcommand{\Z}{\mathbb{Z}}
\newcommand{\F}{\mathbb{F}}
\newcommand{\C}{\mathbb{C}}
\newcommand{\Q}{\mathbb{Q}}
\newcommand{\R}{\mathbb{R}}
\newcommand{\del}{\delta}
\newcommand{\om}{\omega}
\newcommand{\Qscr}{\mathscr{Q}}
\renewcommand{\P}{\mathbb{P}}
\newcommand{\mfp}{\mathfrak{p}}
\newcommand{\Gal}{\operatorname{Gal}}
\newcommand{\Vol}{\operatorname{Vol}}
\newcommand{\Aut}{\operatorname{Aut}}
\newcommand{\SO}{\operatorname{SO}}
\newcommand{\Tr}{\operatorname{Tr} }
\newcommand{\Frob}{\operatorname{Frob}}
\newcommand{\Sel}{\operatorname{Sel}}
\newcommand{\Hilb}{\operatorname{Hilb}}
\newcommand{\Det}{\operatorname{Det}}
\newcommand{\sha}{\operatorname{Sha}}

\usepackage{enumerate}
\begin{document}
% 8.5in paper width -2x1in margin = 6.5in text width
\setlength{\topmargin}{-2mm}



% 11in paper height -2x1in margin = 9in text height


\begin{center}{\bf Homework 4, Math 701 -- Frank Thorne (thorne@math.sc.edu)}
\end{center}
\begin{center}
{\bf Instructions:} You are welcome and encouraged to collaborate, but please write up your own solutions. \\

Due  Friday, November 20, 2026.
\end{center}

\begin{enumerate}

\item
Let $F$ be a field. Prove that the matrix ring $M_n(F)$ is {\itshape simple}: it has no two-sided ideals other than $0$ and itself.

\item
Prove (with an explicit description) that $M_2(\R)$ contains a subring isomorphic to $\C$.

\item
Use the Chinese Remainder Theorem to prove the existence of a partial fraction decomposition: If 
$g(x) = \prod_{i=1}^n (x - a_i) \in \C[x]$ is a polynomial of degree $n$ with distinct roots, and $f(x)$ is any polynomial of lesser degree,
then we may write
\[
\frac{f(x)}{g(x)} = \frac{ b_1 }{x - a_1} + \cdots + \frac{b_n}{x - a_n}
\]
for complex numbers $b_i$.

\item
Prove:
\begin{enumerate}
\item
The maximal ideals of $\C[x]$ are $(x - a)$ for $a \in \C$; 
\item
The only prime, nonmaximal ideal of $\C[x]$ is $(0)$. (Recall that every maximal ideal is prime.)
\end{enumerate}

\item 
Let $R$ be the ring of continuous functions from $[0, 1]$ to $\mathbb{R}$. 
\begin{enumerate}
\item
Determine the units of $R$. 
\item 
Determine all the zero divisors of $R$.
\item 
For each $0 \leq t \leq 1$, prove that $\{ f \in R: f(t) = 0 \}$ is a maximal ideal.
\item 
Prove that every maximal ideal is of this form. 

{\itshape Hint:} You will probably need a bit of real analysis in your proof. 
\end{enumerate}

\end{enumerate}


\end{document}


\item
({\bf Row Rank Equals Column Rank})
The {\bf row rank} and {\bf column rank} of an $n \times m$ 
matrix $A$ are the dimensions of the subspaces of $F^n$
and $F^m$, respectively, spanned by its rows and columns respectively.

Prove\footnote{Credit: I learned of this proof from Marc van Leeuwen} that the row
rank and column rank are equal as follows:
\begin{itemize}
\item
Show that the column rank is the smallest number $r$ for which the homomorphism
$\phi$ corresponding to $A$ can be written as a composition of two linear maps
$F^m \rightarrow F^r \rightarrow F^n$.
\item
Conclude that $r$ is the smallest number for which one can write $A = BC$,
with $B$ an $n \times r$ and $C$ an $r \times m$ matrix respectively.
\item
Taking transposes, conclude the result.
\end{itemize}
\item
Let $V$ be a vector space, possibly infinite dimensional.
\begin{enumerate}
\item
Prove that the map
\[
v \rightarrow \{ f \rightarrow f(v) \}
\]
defines an injective map from $V$ to the double dual $V^{**} = (V^*)^*$.
\item
Prove that this map is not surjective if $V$ is infinite dimensional.
\end{enumerate}

\item
Let $V$ be a finite dimensional vector space, and let $A$ and $B$ be the matrices of a linear transformation $\phi$
with respect to two different bases.

Prove that $B = M A M^{-1}$ for some matrix $M$, and describe $M$ explicitly.

\item
Exhibit a $2 \times 2$ matrix with real entries which does not have any eigenvalues over $\mathbb{R}$. (Better still:
can you characterize all such matrices?)

\end{enumerate}

\end{document}
